\subsection{NASA \textit{Autonomous NanoTechnology Swarm} (ANTS)}

O projeto \textit{Autonomous NanoTechnology Swarm} (ANTS) foi uma iniciativa da NASA desenvolvida no Centro de Voo Espacial Goddard, em colaboração com a Universidade de Ulster e a \textit{Science Applications International Corporation} (SAIC). O projeto partia de uma premissa motivada tanto por limitações físicas quanto operacionais: nas missões de exploração do espaço profundo, o atraso de comunicação entre uma sonda e o controle em Terra pode chegar a dezenas de minutos em cada sentido, tornando inviável qualquer forma de supervisão humana em tempo real. As missões anteriores, como a DS1 (\textit{Deep Space 1}) e a Mars Pathfinder, já haviam explorado graus limitados de autonomia embarcada; o ANTS representou um salto conceitual além desses precedentes, propondo um sistema inteiramente distribuído e sem controle centralizado em Terra \cite{W1999392360, truszkowski2004}.

O conceito central do ANTS era uma frota de mil nanoespaçonaves da classe pico --- cada uma com massa inferior a um quilograma --- encarregadas de prospectar o cinturão de asteroides em busca de recursos minerais e compostos com relevância astrobiológica. Propulsionadas por velas solares, as espaçonaves seriam lançadas a partir de uma órbita no ponto de Lagrange L1 do sistema Terra-Lua. A frota organizava-se em três papéis funcionais complementares: \textit{trabalhadores} (aproximadamente 80\% da frota), responsáveis pela coleta de dados científicos com instrumentos especializados como magnetômetros e espectrômetros; \textit{coordenadores}, encarregados de definir prioridades de missão e organizar o esforço coletivo; e \textit{mensageiros}, responsáveis pela comunicação entre os diferentes subgrupos e com a estação terrestre \cite{truszkowski2004}. O projeto previa dez tipos de especialistas distribuídos em dez subgrupos de aproximadamente cem espaçonaves cada.

Do ponto de vista da computação autônoma, o ANTS materializou todas as quatro propriedades Self-X descritas na Seção \ref{sec:propriedades}. A \textbf{auto-cura} era a propriedade mais crítica: dado que entre 60\% e 70\% das espaçonaves eram esperadas a ser perdidas ao entrar no cinturão de asteroides, as sobreviventes precisavam reorganizar-se autonomamente para redistribuir papéis e cobrir as lacunas deixadas pelas naves perdidas \cite{W1999392360}. A \textbf{auto-configuração} manifestava-se na capacidade de cada espaçonave de determinar sua posição na frota, estabelecer rotas de comunicação com os vizinhos disponíveis e assumir um papel diferente quando necessário. A \textbf{auto-otimização} operava na alocação dinâmica de recursos da missão --- decidindo quais asteroides priorizar com base nos dados já coletados pelos trabalhadores. Já a \textbf{auto-proteção} incluía mecanismos para evitar colisões entre naves do enxame e gerenciar autonomamente o consumo de energia das velas solares.

O modelo de controle do ANTS pode ser interpretado como uma implementação distribuída do laço MAPE-K (Seção \ref{sec:mapek}): cada espaçonave executava localmente seu próprio ciclo de monitoramento-análise-planejamento-execução, enquanto o comportamento coletivo emergente do enxame realizava um segundo nível de auto-gerenciamento de escopo global. Esse caráter hierárquico e emergente distinguia o ANTS dos projetos precursores, que dependiam de agentes individuais com comportamento predefinido. O desafio central residia em garantir que comportamentos corretos no nível local produzissem resultados coerentes e úteis no nível global --- um problema que aproximava o projeto da pesquisa em sistemas multiagentes e inteligência de enxame \cite{truszkowski2004}.

\subsection{IBM \textit{Project eLiza}: Servidores Auto-Gerenciáveis}

A publicação do manifesto da IBM em outubro de 2001 levantou uma questão imediata: como traduzir os princípios de auto-gerenciamento em produtos comerciais concretos? A resposta foi o \textit{Project eLiza}, iniciativa lançada como o maior foco individual de investimento na divisão de servidores da empresa à época \cite{ganek2003}. O projeto estabelecia a linha \textit{eServer} da IBM --- que incluía as famílias zSeries (mainframes), pSeries, iSeries e xSeries --- como plataforma de referência para a implementação progressiva das propriedades Self-X. A meta declarada era transformar a infraestrutura de TI corporativa em um sistema autônomo capaz de gerenciar a si próprio, reduzindo drasticamente a carga de administração humana em ambientes de alta complexidade \cite{ibm2001, ganek2003}.

O \textit{Project eLiza} introduziu capacidades autônomas concretas nos servidores zSeries, que já constituíam a espinha dorsal de grande parte das operações bancárias, de telecomunicações e governamentais globais. A \textbf{auto-configuração} foi implementada por meio de mecanismos que permitiam ao hardware e ao \textit{middleware} reconfigurar subsistemas e recursos de forma autônoma tanto na inicialização quanto em tempo de execução, sem intervenção do operador \cite{ganek2003}. A \textbf{auto-otimização} materializou-se no \textit{Intelligent Resource Director} (IRD), componente do z/OS que monitorava continuamente a utilização de recursos e redistribuía cargas de trabalho entre partições lógicas (\textit{LPARs}) para maximizar o \textit{throughput} e o cumprimento dos acordos de nível de serviço \cite{ganek2003}.

Para a \textbf{auto-cura}, o projeto introduziu o \textit{System Automation for OS/390}, uma camada de política baseada em regras capaz de detectar falhas em aplicações e recursos de sistema e acionar recuperações automáticas --- iniciando, parando, monitorando e recuperando serviços z/OS sem intervenção humana \cite{ganek2003}. Complementarmente, o \textit{msys for Operations} automatizava tarefas operacionais rotineiras, reduzindo a dependência de procedimentos manuais suscetíveis a erro humano. A \textbf{auto-proteção} integrava mecanismos de detecção de intrusão e isolamento de falhas diretamente ao sistema operacional, garantindo a integridade de aplicações críticas diante de ameaças externas ou degradação interna de componentes.

O \textit{Project eLiza} foi simultaneamente um produto comercial e uma prova de conceito científica. A IBM dedicou o volume 42 do \textit{IBM Systems Journal}, publicado em 2003, integralmente ao tema da computação autônoma, reunindo artigos de pesquisadores da empresa sobre os fundamentos teóricos e as implementações práticas do projeto \cite{ganek2003}. Ao demonstrar que as propriedades Self-X podiam ser incorporadas em sistemas de missão crítica em produção, o \textit{Project eLiza} estabeleceu que a computação autônoma não era apenas uma visão acadêmica, mas uma abordagem de engenharia viável para infraestruturas reais --- pavimentando o caminho para a gestão autônoma de larga escala que viria a caracterizar a era subsequente da computação em nuvem.
